Spring:
-------
1.why spring?
        ->it is a light waited process used to build enterprice applications
        ->IT Supports 
                ->Dependeny Injectioon
                -> Inversion of Control
                ->Web application development
                ->Database integration
                -> Security
2.Before Spring
        ->Objected is created manually 
        ->classes where tigtly coupled(one class directly depends on a specific implementation of another class.)
        -> Testing was difficult
        -> Database and transaction handling required more code

3.Advantages of spring
        ->Reduce code
        ->Create loosely coupled applications
        ->testing is easy
        1. Lightweight
        2. Loosely Coupled
        3. Dependency Injection
        4. Inversion of Control
        5. Easy Testing
        6. Transaction Management
        7. Database Integration
        8. Web Application Development
        9. Security
        10. Supports Modular Development


4.Dependecy Injection:
         providing the needed thing outside the class than creating it inside the class is know as depenedency injection
         Dependency Injection is the process of providing the required dependency of a class from outside instead of creating it inside the class.
5.Spring Container:
        creating Object
        Managet the objecs
        Injecting dependency
        manageing bean life cycle
        
Types of spring Container:
------------------------
        1.Spring Factory
        2.Application Context
===================================================================================================================================================
Bean
----
        bean is an object managed by spring container
        
        Example:
        
        @Service
        class EmployeeService {
        
        }
        Spring creates an object of EmployeeService and manages it.

        That object is called a Spring Bean.

        Class
          ↓
        @Service
          ↓
        Spring Container
          ↓
        Object created
          ↓
        Spring Bean

Component Annotation
        1.@service->it is used to write business logics inside the class
        2.@repository->Handels database request and response
        3.@controller->use to handel web request and response 
        4.@Component tells Spring that the class should be managed by the Spring Container.
        5.@RestController->It is used to create rest API's 





































        
Spring Boot:
--------------
        Spring boot is an Frame Work for Java that is build on top of Spring that simplifies Application Development
        It Supports
                ->web Application
                ->Rest Api's
                ->microservices
                ->security
                ->and for cloud Deployment


why Sprint boot?
            ->Fast Development
            ->Auto-configuration
            ->microservices friendly
            ->Seamless Integration

Features of SprintBoot:
-----------------------
                ->Auto-Configuration: Automatically configures components (like Hibernate, JPA) based on dependencies. 
                ->Easy Maintenance and Creation of REST Endpoints: With annotations like @RestController, @GetMapping and @PostMapping, creating REST endpoints is straightforward.


Annotations Explanation:

1.Core Basic Annotation:

@SpringBootApplication:
----------------------
                  ->Main Entry Point of Spring Boot
                  ->example:
                  @SpringBootApplication
                              public class DemoApplication {
                                          public static void main(String[] args) {
                                          SpringApplication.run(DemoApplication.class, args);
                                                 }
                                          }

@Component
----------
                  -It tells spring This is a class manage it for manage
                  @Component
                  public class EmailService {}

                        
@Autowired
----------

      ->  @autowired:@Autowired is used to automatically connect (inject) one class into another in Spring Boot.
      ->  we usually connect Repository / service to the controller

            example:
              @Autowired
             private bookrepo bookRepo;

@Bean:
-----
      ->it is used to create Object Manually and gives it to Spring
      @Bean
            public ModelMapper modelMapper()
            {
                  return new ModelMapper();
            }
        

2.Steriotype Annotation:

@controller
------------
            ->handels web pages 
                @Controller
                public class PageController {
                }

@service
---------
            ->it Holds Business logics
            
=======
Spring Boot:
--------------
        Spring boot is an Frame Work for Java that is build on top of Spring that simplifies Application Development
        It Supports
                ->web Application
                ->Rest Api's
                ->microservices
                ->security
                ->and for cloud Deployment


why Sprint boot?
            ->Fast Development
            ->Auto-configuration
            ->microservices friendly
            ->Seamless Integration

Features of SprintBoot:
-----------------------
                ->Auto-Configuration: Automatically configures components (like Hibernate, JPA) based on dependencies. 
                ->Easy Maintenance and Creation of REST Endpoints: With annotations like @RestController, @GetMapping and @PostMapping, creating REST endpoints is straightforward.


Annotations Explanation:

1.Core Basic Annotation:

@SpringBootApplication:
----------------------
                  ->Main Entry Point of Spring Boot
                  ->example:
                  @SpringBootApplication
                              public class DemoApplication {
                                          public static void main(String[] args) {
                                          SpringApplication.run(DemoApplication.class, args);
                                                 }
                                          }

@Component
----------
                  -It tells spring This is a class manage it for manage
                  @Component
                  public class EmailService {}

                        
@Autowired
----------

      ->  @autowired:@Autowired is used to automatically connect (inject) one class into another in Spring Boot.
      ->  we usually connect Repository / service to the controller

            example:
              @Autowired
             private bookrepo bookRepo;

@Bean:
-----
      ->it is used to create Object Manually and gives it to Spring
      @Bean
            public ModelMapper modelMapper()
            {
                  return new ModelMapper();
            }
        

2.Steriotype Annotation:

@controller
------------
            ->handels web pages 
                @Controller
                public class PageController {
                }

@service
---------
            ->it Holds Business logics


@Reposiory
----------
            ->it is a database Layer


3.Request Mappinng Annotations 

@RequestMapping
            
>>>>>>> Stashed changes
